\begin{figure}[H]
\centering
\begin{tikzpicture}[font=\small,>=Stealth]

\draw[very thick] (0,0)--(6,0) node[right] {$K=[0,1]$};
\foreach \x/\y in {0/.4,1/.65,2/.9,3/1.15,4/1.4,5/1.65,6/1.9}{
\draw[gray,dashed] (\x-1.15,\y)--(\x+1.15,\y);}
\foreach \x/\y in {0/.4,2/.9,4/1.4,6/1.9}{
\draw[blue,very thick] (\x-1.15,\y)--(\x+1.15,\y);
\draw[blue,fill=white] (\x-1.15,\y) circle (1.5pt) (\x+1.15,\y) circle (1.5pt);}
\node[align=center] at (3,-.65) {blue intervals are retained members; gray intervals may be discarded};

\end{tikzpicture}
\caption{A finite subcover retains members of the original cover. The vertical offsets separate the intervals visually; all intervals live on the same horizontal line. This schematic illustrates one cover, whereas compactness controls every open cover.}
\label{fig:compact-subcover}
\end{figure}
